\cotefile{57}
\scanpage{3}
\uline{Local Picard schemes}.

Let $A$ be a complete local ring, \deleted{\unsure{X}} $S=\operatorname{Spec}(A)$, $U=S-\handwritten{\{}s\handwritten{\}}$ the complement of the closed point of $S$. We wish to interprete the Picard \deleted{\unsure{gro}} group $\operatorname{Pic}(U)$ as \deleted{some} the group $P(k)$ of $k$-valued points of a group-scheme $P$ defined over the residue field $k$ of $A$. Assume for simplicity that $U$ is regular, and that there exists a resolution of singularities for $S$ not changing $U$, i.e. a proper morphism
\[
f:X\handwritten{\longrightarrow}S
\]
inducing an isomorphism\deleted{\unsure{s}} $f^{-1}(U)\handwritten{\simeq}U$, such that $X$ be regular and equal to the closure of $f^{-1}(U)$ (this open set will be identified with $U$). Let $X_o$ be the special fiber of $X$, so that
\[
U=X-X_o\ .
\]
Because $X$ is regular, we have an exact sequence
\[
D\handwritten{\longrightarrow}\operatorname{Pic}(X)\handwritten{\longrightarrow}\operatorname{Pic}(U)\handwritten{\longrightarrow}0\ ,
\tag*{\handwritten{(1)}}
\]
where $D$ is the group of divisors on $X$ concentrated on $X_o$ (which is a free group of finite type $\underline{Z}^{I}$, $I$ being the set of irreducible components of $X_o$). On the other hand, by the existence theorem for coherent sheaves, it is known that
\[
\operatorname{Pic}(X)\handwritten{\simeq}\underset{n}{\underset{\handwritten{\longleftarrow}}{\lim}}\operatorname{Pic}(X_n)\ ,
\tag*{\handwritten{(2)}}
\]
where $X_n=X\mathrm{x}_S S_n$, $S_n=\operatorname{Spec}(A/\underline{m}^{n+1})$. Assume we can find group schemes $P_n$ over $k$ and isomorphisms
\[
\operatorname{Pic}(X_n)\handwritten{\xrightarrow{\sim}}P_n(k)\ ,
\tag*{\handwritten{(3)}}
\]
and transition morphisms $P_m\handwritten{\longrightarrow}P_n$ $(m\handwritten{\geq}n)$ compatible with the transition morphisms for the $\operatorname{Pic}(X_n)$. Denoting by $P$ the pro-group $(P_n)$, or if one prefers, in case the transition morphisms $P_m\handwritten{\longrightarrow}P_n$ are affine, the inverse limit (EGA IV 8) of the $P_n$, we get a pro-group scheme resp. a group-scheme $P$, and an isomorphism
% archive sheet number: 2
